
\documentclass{article}
\usepackage{colortbl}
\usepackage{makecell}
\usepackage{multirow}
\usepackage{supertabular}

\begin{document}

\newcounter{utterance}

\centering \large Interaction Transcript for game `bbh', experiment `formal\_fallacies', episode 158 with reflexion{-}Qwen3.5{-}4B.
\vspace{24pt}

{ \footnotesize  \setcounter{utterance}{1}
\setlength{\tabcolsep}{0pt}
\begin{supertabular}{c@{$\;$}|p{.15\linewidth}@{}p{.15\linewidth}p{.15\linewidth}p{.15\linewidth}p{.15\linewidth}p{.15\linewidth}}
    \# & \multicolumn{2}{c}{Player} && \multicolumn{2}{c}{Game Master} \\
    \hline

    \theutterance \stepcounter{utterance}  
    & & & \multicolumn{4}{p{0.6\linewidth}}{
        \cellcolor[rgb]{0.9,0.9,0.9}{
            \makecell[{{p{\linewidth}}}]{
                \texttt{\tiny{[P1$\langle$GM]}}
                \texttt{You have attempted this same task before and did not succeed. These are the plans you wrote after those attempts.} \\
\texttt{Trial \#1: New plan: I failed because I misinterpreted the logical structure of Premise 3. I treated "Being a close friend of Shanika is necessary for being a close friend of Francine" as \$F \textbackslash{}rightarrow S\$, but the correct translation of "X is necessary for Y" is \$Y \textbackslash{}rightarrow X\$ (i.e., \$F \textbackslash{}rightarrow S\$ is actually correct in my notation, but I failed to realize that Premise 2 (\$F \textbackslash{}rightarrow C\$) combined with Premise 1 (\$F \textbackslash{}rightarrow C \textbackslash{}lor A \textbackslash{}lor S\$) makes the disjunction redundant and provides no information about \$A\$). Specifically, knowing \$F\$ implies \$C\$ makes the condition \$(C \textbackslash{}lor A \textbackslash{}lor S)\$ always true for any \$F\$, rendering \$A\$ irrelevant to the truth of \$F\$. I should have immediately recognized that Premise 2 alone is sufficient to satisfy Premise 1 for any \$F\$, meaning \$A\$ is never forced to be true. Next time, I will strictly translate "necessary for" as "if consequent then antecedent" and check if any premise makes the disjunction in another premise redundant before concluding validity.} \\
\\ 
\\ 
\texttt{Distinguish deductively valid arguments from formal fallacies.} \\
\\ 
\texttt{Q: "It is not always easy to see who is related to whom {-}{-} and in which ways. The following argument pertains to this question: To begin with, Lesley is a close friend of Fernando. Moreover, being a close friend of Fernando or a schoolmate of Lowell is sufficient for being a great{-}grandfather of Leroy. It follows that Lesley is a great{-}grandfather of Leroy."} \\
\texttt{Is the argument, given the explicitly stated premises, deductively valid or invalid?} \\
\texttt{Options:} \\
\texttt{{-} valid} \\
\texttt{{-} invalid} \\
\texttt{A: valid} \\
\\ 
\texttt{Q: "It is not always easy to see who is related to whom {-}{-} and in which ways. The following argument pertains to this question: Whoever is not a great{-}grandfather of Clyde is a stepbrother of Brian. Being an ancestor of Dana is sufficient for not being a great{-}grandfather of Clyde. We may conclude: Everyone who is an ancestor of Dana is a stepbrother of Brian, too."} \\
\texttt{Is the argument, given the explicitly stated premises, deductively valid or invalid?} \\
\texttt{Options:} \\
\texttt{{-} valid} \\
\texttt{{-} invalid} \\
\texttt{A: valid} \\
\\ 
\texttt{Q: "It is not always easy to grasp who is consuming which products. The following argument pertains to this question: Every infrequent user of Paul Mitchell shampoo is either a rare consumer of Nioxin shampoo or a loyal buyer of Caress soap, or both. No regular consumer of Lush soap is a rare consumer of Nioxin shampoo and, in the same time, a loyal buyer of Caress soap. It follows that whoever is an infrequent user of Paul Mitchell shampoo is not a regular consumer of Lush soap."} \\
\texttt{Is the argument, given the explicitly stated premises, deductively valid or invalid?} \\
\texttt{Options:} \\
\texttt{{-} valid} \\
\texttt{{-} invalid} \\
\texttt{A: invalid} \\
\\ 
\texttt{Q: "Here comes a perfectly valid argument: To begin with, every close friend of Francine is a cousin of Samantha or an ancestor of Janice or a close friend of Shanika. Moreover, every close friend of Francine is a cousin of Samantha. Finally, being a close friend of Shanika is necessary for being a close friend of Francine. We may conclude: Being an ancestor of Janice is necessary for being a close friend of Francine."} \\
\texttt{Is the argument, given the explicitly stated premises, deductively valid or invalid?} \\
\texttt{Options:} \\
\texttt{{-} valid } \\
\texttt{{-} invalid} \\
\texttt{A:} \\
            }
        }
    }
     \\ \\

    \theutterance \stepcounter{utterance}  
    & \multicolumn{4}{p{0.6\linewidth}}{
        \cellcolor[rgb]{0.561,0.737,0.561}{
            \makecell[{{p{\linewidth}}}]{
                \texttt{\tiny{[P1$\rangle$GM]}}
                \texttt{invalid} \\
\\ 
\texttt{**Reasoning:**} \\
\texttt{Let's break down the argument using propositional logic:} \\
\texttt{*   Let \$F\$ = Being a close friend of Francine.} \\
\texttt{*   Let \$S\$ = Being a cousin of Samantha.} \\
\texttt{*   Let \$A\$ = Being an ancestor of Janice.} \\
\texttt{*   Let \$Sh\$ = Being a close friend of Shanika.} \\
\\ 
\texttt{**Premises:**} \\
\texttt{1.  \$F \textbackslash{}rightarrow (S \textbackslash{}lor A \textbackslash{}lor Sh)\$ (Every close friend of Francine is a cousin of Samantha or an ancestor of Janice or a close friend of Shanika).} \\
\texttt{2.  \$F \textbackslash{}rightarrow S\$ (Every close friend of Francine is a cousin of Samantha).} \\
\texttt{3.  \$Sh \textbackslash{}rightarrow F\$ (Being a close friend of Shanika is necessary for being a close friend of Francine. Note: "X is necessary for Y" translates to \$Y \textbackslash{}rightarrow X\$. Here, \$Sh\$ is necessary for \$F\$, so \$F \textbackslash{}rightarrow Sh\$? Wait, let's re{-}read carefully. "Being a close friend of Shanika is necessary for being a close friend of Francine." This means if you are a friend of Francine, you must be a friend of Shanika. So \$F \textbackslash{}rightarrow Sh\$.} \\
\texttt{    *   *Correction on Trial \#1 analysis*: The user's trial note says they treated it as \$F \textbackslash{}rightarrow S\$ but realized the disjunction was redundant. Let's re{-}evaluate the specific logic of the conclusion.} \\
\texttt{    *   Premise 3: "Being a close friend of Shanika is necessary for being a close friend of Francine."} \\
\texttt{        *   Structure: \$F \textbackslash{}rightarrow Sh\$. (If \$F\$, then \$Sh\$).} \\
\\ 
\texttt{**Conclusion:**} \\
\texttt{*   \$A \textbackslash{}rightarrow F\$ (Being an ancestor of Janice is necessary for being a close friend of Francine).} \\
\texttt{    *   Structure: \$F \textbackslash{}rightarrow A\$. (If \$F\$, then \$A\$).} \\
\\ 
\texttt{**Analysis:**} \\
\texttt{1.  From Premise 2 (\$F \textbackslash{}rightarrow S\$) and Premise 1 (\$F \textbackslash{}rightarrow (S \textbackslash{}lor A \textbackslash{}lor Sh)\$):} \\
\texttt{    *   If \$F\$ is true, then \$S\$ is true.} \\
\texttt{    *   If \$S\$ is true, then \$(S \textbackslash{}lor A \textbackslash{}lor Sh)\$ is automatically true.} \\
\texttt{    *   Therefore, Premise 1 is redundant given Premise 2. It does not add any new constraints.} \\
\texttt{2.  From Premise 3 (\$F \textbackslash{}rightarrow Sh\$):} \\
\texttt{    *   If \$F\$ is true, then \$Sh\$ is true.} \\
\texttt{3.  Does this imply the Conclusion (\$F \textbackslash{}rightarrow A\$)?} \\
\texttt{    *   We know that if \$F\$ is true, then \$S\$ is true AND \$} \\
            }
        }
    }
    & & \\ \\

    \theutterance \stepcounter{utterance}  
    & & & \multicolumn{2}{p{0.3\linewidth}}{
        \cellcolor[rgb]{0.9,0.9,0.9}{
            \makecell[{{p{\linewidth}}}]{
                \texttt{\tiny{[GM$|$GM]}}
                \texttt{Target: invalid} \\
            }
        }
    }
    & & \\ \\

    \theutterance \stepcounter{utterance}  
    & & & \multicolumn{2}{p{0.3\linewidth}}{
        \cellcolor[rgb]{0.9,0.9,0.9}{
            \makecell[{{p{\linewidth}}}]{
                \texttt{\tiny{[GM$|$GM]}}
                \texttt{game\_result = LOSE} \\
            }
        }
    }
    & & \\ \\

    \theutterance \stepcounter{utterance}  
    & & \multicolumn{4}{p{0.6\linewidth}}{
        \cellcolor[rgb]{0.95,0.95,0.95}{
            \makecell[{{p{\linewidth}}}]{
                \texttt{\tiny{[P1$\leftrightarrow$P]}}
                \texttt{Post{-}episode Reflexion (not sent to GM):} \\
\texttt{New plan: I failed because I incorrectly translated the logical structure of Premise 3. I treated "Being a close friend of Shanika is necessary for being a close friend of Francine" as implying \$F \textbackslash{}rightarrow A\$ (or confused the direction of implication), whereas "X is necessary for Y" means \$Y \textbackslash{}rightarrow X\$. In this case, it means \$F \textbackslash{}rightarrow Sh\$. Since Premise 2 states \$F \textbackslash{}rightarrow S\$, and Premise 1 is a disjunction \$(S \textbackslash{}lor A \textbackslash{}lor Sh)\$, knowing \$F\$ implies \$S\$ (which satisfies the disjunction) and \$Sh\$ (which also satisfies it). However, this provides **no information** about \$A\$. The conclusion claims \$F \textbackslash{}rightarrow A\$, but the premises only guarantee \$F \textbackslash{}rightarrow S\$ and \$F \textbackslash{}rightarrow Sh\$. \$A\$ is completely unconstrained by the premises. Next time, I will strictly apply the rule: "X is necessary for Y" \$\textbackslash{}equiv\$ \$Y \textbackslash{}rightarrow X\$. I will then check if the conclusion's antecedent/consequent relationship is actually supported by the premises. If the conclusion requires a variable (\$A\$) that is not forced to be true by the premises (even if other variables in a disjunction are true), the argument is invalid. I will explicitly verify that no premise forces \$A\$ to be true whenever \$F\$ is true.} \\
            }
        }
    }
    & \\ \\

\end{supertabular}
}

\end{document}
